\documentclass[UTF8,fontset=fandol]{ctexbeamer}

\usetheme{tiev2026}

\usepackage{booktabs}
\usepackage{graphicx}
\usepackage{amsmath}
\usepackage[numbers]{natbib}

\title{汇报题目 / Presentation Title}
\subtitle{副标题或项目名称}
\author[姓名]{姓名 同济智能无人系统实验室}
\date{2026 年 5 月}

\begin{document}

% 版式 1：标题页
\maketitle

% 版式 2：标题和内容
\begin{frame}[t,fragile]{标题和内容}
  \begin{itemize}
    \item 这里放置本页的核心结论，尽量使用短句\cite{latexcompanion}。
    \item 普通内容页复刻 PPT 中的左侧蓝色竖条、右上角校徽与简化 logo。
    \item 页脚文字和页码由主题自动生成，可在导言区重定义 \verb|\TIEVFooterText|。
  \end{itemize}

  \vspace{0.35cm}
  \begin{block}{重点信息}
    可以用 block 承载方法、实验设置或阶段性结论。
  \end{block}
\end{frame}

% 版式 3：双栏内容
\begin{frame}[t]{双栏版式}
  \begin{columns}[T,onlytextwidth]
    \begin{column}{0.48\textwidth}
      \begin{block}{左侧内容}
        \begin{itemize}
          \item 研究背景
          \item 方法动机
          \item 关键假设
        \end{itemize}
      \end{block}
    \end{column}
    \begin{column}{0.48\textwidth}
      \begin{block}{右侧内容}
        \begin{enumerate}
          \item 模块一
          \item 模块二
          \item 模块三
        \end{enumerate}
      \end{block}
    \end{column}
  \end{columns}
\end{frame}

% 版式 4：章节或大标题页
\TIEVSectionPage{章节页 / Section Divider}

% 版式 5：空白或大图页
\begin{TIEVBlankFrame}
  \vspace*{0.70cm}
  \begin{center}
    {\fontsize{22}{28}\selectfont\bfseries 空白或大图版式\par}
    \vspace{0.45cm}
    \fbox{\rule{0pt}{4.15cm}\rule{10.6cm}{0pt}}
  \end{center}
\end{TIEVBlankFrame}

% 引用与参考文献
\begin{frame}[t,fragile]{引用示例}
  使用 \verb|\citep| 可以生成括号引用，例如 \citep{latexcompanion}；
  使用 \verb|\citet| 可以生成文本引用，例如 \citet{latexcompanion}。

  \vspace{0.35cm}
  将 BibTeX 条目写入 \verb|references.bib|，编译时会自动生成参考文献。
\end{frame}

\begin{frame}[allowframebreaks]{参考文献}
  \bibliographystyle{plainnat}
  \bibliography{references}
\end{frame}

% 结束页：PPT 中的“敬请批评指正”
\TIEVThanksPage[姓名 \quad email@tongji.edu.cn]{敬请批评指正！}

\end{document}
